\chapter{Bevezetés}
\label{ch:intro}

Egy növényzet digitálisan való elkészítése egy rendkívül fontos témakör az informatika területén is. A számítógépes játék iparban problémakör a hatalmas területek lefedése különféle növényfajokkal úgy, hogy az esztétikailag vonzó legyen. Szimulációs applikációk esetén pedig a megjelenített növényzettől sokszor elvárt kritérium annak realisztikusabb kinézete, felépítése.  Ennek megoldására rengeteg procedurálisan dolgozó eszköz és technika létezik, amelyek az adott célokra jobban vagy rosszabbul teljesítenek.
\newline
A témakör keretein belül készített szoftver is egy ilyen procedurális vegetációt produkáló applikáció, amellyel úgynevezett tájékozódási futó térképek segítségével képes a felhasználó különböző növényzetek leírását generáltatni. A tájfutótérképek elég részletesen adnak ugyanis információt egy területen elhelyezkedő növényzetről és egyéb, például annak terjedését akadályozó objektumokról. Ezeket az adatokat értelmezve és összevetve pedig olyan területeket képes a program meghatározni, amelyeken különböző növények tudnak megjelenni.
\newline
Ezen eszközök alkalmazása digitális világok leírására továbbá nem egy új vagy meglepő ötlet, hiszen ezzel a megoldással a már meglévő sztenderdeket és elkészült térképeket tudjuk egy más kontextusban újrafelhasználni. Viszont a sztenderdekk és technológiák frissülése miatt fontosnak véltem ezen témakör alaposabb vizsgálatát.
\newline 
A következőkben megtalálható a felhasználói dokumentációban, amiben többek között a különböző állítható paraméterek, a felhasználás lépései és a program alapvető korlátai lesznek részletezve. Továbbá a fejlesztői dokumentációban lehet olvasni a különböző dizájn és implementációs döntésekről, a program felépítéséről és a limitációk fejlesztői oldaláról.  

